\subsection{window-\texorpdfstring{$6$}{6} 的 Coxeter 三项链、Fourier 选择律与长根张量正规形}\label{subsec:conclusion-window6-coxeter-necklace-fourier-doublet-normal-form}
在 window-\(6\) 的 root--Cartan 字典、极小 Lie 完成、纤维多重度场的坐标钉扎以及 boundary 三轴的中心解释都已固定之后，还可进一步把 cyclic 根层内部的 Coxeter 传输、多重度场的 orbitwise Fourier 选择律与 long-sector 的内部张量对称压缩为同一条终端链。其核心现象是：表 \ref{tab:fold6_b3c3_root_dictionary_b3} 上的 \(B_3\) Coxeter 元把 \(X_6^{\mathrm{cyc}}\) 重新组织为一条短根六项链与两条长根六项链；其中两条长根项链承载完全相同的 bare fold profile \((4,2,3,4,2,3)\)，从而强制出现精确的内部 \(GL_2\) 对称与一个 qutrit--sign--doublet 张量正规形；相反，短根层的 projective 对称只在 \(\{\pm e_1\}\) 上发生唯一破缺。

\begin{theorem}[window-\texorpdfstring{$6$}{6} 的 Coxeter 三项链分解]\label{thm:conclusion-window6-coxeter-three-necklace-decomposition}
在 \(B_3\) 的标准根空间 \(\RR^3\) 中取简单根
\[
\alpha_1=e_1-e_2,\qquad \alpha_2=e_2-e_3,\qquad \alpha_3=e_3,
\]
并令
\[
c:=s_{\alpha_1}s_{\alpha_2}s_{\alpha_3}\in W(B_3).
\]
则
\[
c(x_1,x_2,x_3)=(-x_3,x_1,x_2),\qquad c^6=1,\qquad c^3=-I.
\]
再记 \(\iota_B:X_6^{\mathrm{cyc}}\xrightarrow{\sim}\Phi(B_3)\) 为表 \ref{tab:fold6_b3c3_root_dictionary_b3} 的 \(B_3\) 根字典，并定义搬运到 \(X_6^{\mathrm{cyc}}\) 上的置换
\[
\mathfrak c:=\iota_B^{-1}\circ c\circ \iota_B.
\]
则 \(\mathfrak c\) 为 \(X_6^{\mathrm{cyc}}\) 上的一个阶 \(6\) 置换，并且
\[
X_6^{\mathrm{cyc}}=\mathcal O_{\mathrm{s}}\sqcup \mathcal O_{+}\sqcup \mathcal O_{-}
\]
分解为三个互不相交的 \(6\)-循环，其中
\[
\mathcal O_{\mathrm{s}}
=
\{\texttt{100000},\texttt{010000},\texttt{001000},\texttt{000001},\texttt{000010},\texttt{000100}\},
\]
\[
\mathcal O_{+}
=
\{\texttt{000000},\texttt{010001},\texttt{001010},\texttt{010100},\texttt{010101},\texttt{010010}\},
\]
\[
\mathcal O_{-}
=
\{\texttt{101000},\texttt{001001},\texttt{100010},\texttt{100100},\texttt{000101},\texttt{101010}\}.
\]
其中 \(\mathcal O_{\mathrm{s}}\) 恰对应六条短根，而 \(\mathcal O_{\pm}\) 恰对应长根集合中的两条 Coxeter 轨道。
\end{theorem}
\begin{proof}
直接计算三条简单反射得
\[
s_{\alpha_1}(x_1,x_2,x_3)=(x_2,x_1,x_3),\qquad
s_{\alpha_2}(x_1,x_2,x_3)=(x_1,x_3,x_2),\qquad
s_{\alpha_3}(x_1,x_2,x_3)=(x_1,x_2,-x_3).
\]
故
\[
c(x_1,x_2,x_3)=(-x_3,x_1,x_2).
\]
再直接迭代可得
\[
c^3(x_1,x_2,x_3)=(-x_1,-x_2,-x_3),
\]
即 \(c^3=-I\)，从而 \(c^6=1\)。

对三条种子根分别作用，有
\[
e_1\mapsto e_2\mapsto e_3\mapsto -e_1\mapsto -e_2\mapsto -e_3\mapsto e_1,
\]
\[
e_1+e_2\mapsto e_2+e_3\mapsto -e_1+e_3\mapsto -e_1-e_2\mapsto -e_2-e_3\mapsto e_1-e_3\mapsto e_1+e_2,
\]
\[
e_1-e_2\mapsto e_2-e_3\mapsto e_1+e_3\mapsto -e_1+e_2\mapsto -e_2+e_3\mapsto -e_1-e_3\mapsto e_1-e_2.
\]
将这三条根轨道经表 \ref{tab:fold6_b3c3_root_dictionary_b3} 回读到 \(X_6^{\mathrm{cyc}}\)，便分别得到
\(\mathcal O_{\mathrm{s}},\mathcal O_{+},\mathcal O_{-}\) 的显式词表。
其中 \(\mathcal O_{\mathrm{s}}\) 对应
\[
\{\pm e_1,\pm e_2,\pm e_3\},
\]
即短根集合；其余两条轨道则并成全部 \(12\) 条长根。证毕。
\leanverified{paper\_conclusion\_window6\_coxeter\_three\_necklace\_decomposition}
\end{proof}

\begin{theorem}[Coxeter 作用的谱型与根数重现]\label{thm:conclusion-window6-coxeter-spectrum-root-count-recovery}
记
\[
A_6^{\CC}:=\CC^{X_6},
\]
并以 \(\{\delta_w\}_{w\in X_6}\) 为标准基。令 \(U_{\mathfrak c}\in \operatorname{End}(A_6^{\CC})\) 在 cyclic 扇区上由
\[
U_{\mathfrak c}\delta_w:=\delta_{\mathfrak c(w)}
\qquad (w\in X_6^{\mathrm{cyc}})
\]
给出，并在 boundary 子空间 \(\CC^{X_6^{\mathrm{bdry}}}\) 上取恒等作用。则作为 \(C_6=\langle U_{\mathfrak c}\rangle\)-模，有
\[
A_6^{\CC}\cong \CC[C_6]^{\oplus 3}\oplus \CC^3.
\]
特别地，若 \(\zeta_6=e^{2\pi i/6}\)，则
\[
\operatorname{mult}_{A_6^{\CC}}(1)=6,\qquad
\operatorname{mult}_{A_6^{\CC}}(\zeta_6^k)=3\quad (1\le k\le 5).
\]
因而
\[
\dim\ker(U_{\mathfrak c}-I)=6,\qquad
\dim\ker(U_{\mathfrak c}^3-I)=12.
\]
其中前者恰等于 window-\(6\) 的短根总数，后者恰等于长根总数。
\end{theorem}
\begin{proof}
由定理 \ref{thm:conclusion-window6-coxeter-three-necklace-decomposition}，\(\CC^{X_6^{\mathrm{cyc}}}\) 恰分解为三个长度为 \(6\) 的置换模，每一块都同构于 \(C_6\) 的正则表示 \(\CC[C_6]\)。另一方面，定理 \ref{thm:window6-21-adjoint-weight-multiset} 已把
\[
X_6^{\mathrm{bdry}}=\{\texttt{100001},\texttt{100101},\texttt{101001}\}
\]
识别为三重零权，而推论 \ref{cor:fold-bin-boundary-center-z2} 与定理 \ref{thm:conclusion-window6-boundary-torus-2torsion} 已将它们解释为 boundary 的中心三轴；因此令 \(U_{\mathfrak c}\) 在该三维子空间上取恒等作用，与既有 root--Cartan 解释相容。于是
\[
A_6^{\CC}\cong \CC[C_6]^{\oplus 3}\oplus \CC^3.
\]

正则表示在复数域上分解为六个不同特征标各一次，因此三份正则块对特征值 \(1,\zeta_6,\dots,\zeta_6^5\) 各贡献重数 \(3\)，再加上 boundary 的三维平凡块，便得到
\[
\operatorname{mult}_{A_6^{\CC}}(1)=3+3=6,\qquad
\operatorname{mult}_{A_6^{\CC}}(\zeta_6^k)=3\quad (1\le k\le 5).
\]
于是
\[
\dim\ker(U_{\mathfrak c}-I)=6.
\]
再看 \(U_{\mathfrak c}^3\)。在每个正则块上，只有偶指标特征标满足 \((\zeta_6^k)^3=1\)，故每个正则块对 \(\ker(U_{\mathfrak c}^3-I)\) 贡献 \(3\) 维；再加上 boundary 的 \(3\) 维平凡块，即得
\[
\dim\ker(U_{\mathfrak c}^3-I)=3+3+3+3=12.
\]
最后，由定理 \ref{thm:window6-cyclic-weight-threshold-root-length}，window-\(6\) 的短根总数为 \(6\)，长根总数为 \(12\)，故上面两条维数恰与根数完全一致。证毕。
\end{proof}

\begin{theorem}[window-\texorpdfstring{$6$}{6} 的 Coxeter--fold 块正规形]\label{thm:conclusion-window6-coxeter-fold-block-normal-form}
在定理 \ref{thm:conclusion-window6-coxeter-spectrum-root-count-recovery} 的表示空间上定义纤维多重度对角算子
\[
D_6:A_6^{\CC}\to A_6^{\CC},\qquad
D_6\delta_w:=d_6^{\mathrm{bin}}(w)\,\delta_w.
\]
按定理 \ref{thm:conclusion-window6-coxeter-three-necklace-decomposition} 中
\[
(\mathcal O_{\mathrm{s}};\ \mathcal O_{+};\ \mathcal O_{-};\ \texttt{100001},\texttt{100101},\texttt{101001})
\]
的顺序排列基底，则
\[
U_{\mathfrak c}=P_6\oplus P_6\oplus P_6\oplus I_3,
\]
其中 \(P_6\) 为标准 \(6\)-循环置换矩阵；同时
\[
D_6=D_{\mathrm{s}}\oplus D_{\ell}\oplus D_{\ell}\oplus 2I_3,
\]
其中
\[
D_{\mathrm{s}}=\operatorname{diag}(4,4,4,2,4,4),\qquad
D_{\ell}=\operatorname{diag}(4,2,3,4,2,3).
\]
\end{theorem}
\begin{proof}
\(U_{\mathfrak c}\) 的块形只是定理 \ref{thm:conclusion-window6-coxeter-three-necklace-decomposition} 的矩阵化表述。

对 \(D_6\)，由表 \ref{tab:fold_bin_gauge_m6_fibers} 可直接读出
\[
d_6^{\mathrm{bin}}(\texttt{100000})=
d_6^{\mathrm{bin}}(\texttt{010000})=
d_6^{\mathrm{bin}}(\texttt{001000})=
d_6^{\mathrm{bin}}(\texttt{000010})=
d_6^{\mathrm{bin}}(\texttt{000100})=4,
\]
\[
d_6^{\mathrm{bin}}(\texttt{000001})=2,
\]
故短根轨道 \(\mathcal O_{\mathrm{s}}\) 对应的块正是
\[
D_{\mathrm{s}}=\operatorname{diag}(4,4,4,2,4,4).
\]
同样，
\[
\bigl(
d_6^{\mathrm{bin}}(\texttt{000000}),
d_6^{\mathrm{bin}}(\texttt{010001}),
d_6^{\mathrm{bin}}(\texttt{001010}),
d_6^{\mathrm{bin}}(\texttt{010100}),
d_6^{\mathrm{bin}}(\texttt{010101}),
d_6^{\mathrm{bin}}(\texttt{010010})
\bigr)
=(4,2,3,4,2,3),
\]
\[
\bigl(
d_6^{\mathrm{bin}}(\texttt{101000}),
d_6^{\mathrm{bin}}(\texttt{001001}),
d_6^{\mathrm{bin}}(\texttt{100010}),
d_6^{\mathrm{bin}}(\texttt{100100}),
d_6^{\mathrm{bin}}(\texttt{000101}),
d_6^{\mathrm{bin}}(\texttt{101010})
\bigr)
=(4,2,3,4,2,3),
\]
故两条长根项链的块完全相同，均为 \(D_{\ell}\)。
最后，由定理 \ref{thm:window6-21-adjoint-weight-multiset} 可知 boundary 三词恰为
\[
\texttt{100001},\qquad \texttt{100101},\qquad \texttt{101001},
\]
而表 \ref{tab:fold_bin_gauge_m6_fibers} 表明它们都满足 \(d_6^{\mathrm{bin}}=2\)，故 boundary 块为 \(2I_3\)。证毕。
\leanverified{paper\_conclusion\_window6\_coxeter\_fold\_block\_normal\_form}
\end{proof}

\begin{corollary}[长根射影商与唯一短根定向缺陷]\label{cor:conclusion-window6-long-root-projective-short-root-defect}
记
\[
\Phi_{\mathrm{long}}(B_3):=\{\pm e_i\pm e_j:\ 1\le i<j\le 3\},
\qquad
\Phi_{\mathrm{short}}(B_3):=\{\pm e_1,\pm e_2,\pm e_3\}.
\]
则 \(d_6^{\mathrm{bin}}\) 在长根上满足
\[
d_6^{\mathrm{bin}}(\alpha)=d_6^{\mathrm{bin}}(-\alpha)
\qquad (\alpha\in \Phi_{\mathrm{long}}(B_3)),
\]
因而它下降为射影长根商
\[
\Phi_{\mathrm{long}}(B_3)/\{\pm 1\}
\]
上的函数。相反，在短根上该下降仅失败于唯一一条根线：
\[
d_6^{\mathrm{bin}}(e_1)=4,\qquad d_6^{\mathrm{bin}}(-e_1)=2,
\]
而
\[
d_6^{\mathrm{bin}}(e_2)=d_6^{\mathrm{bin}}(-e_2)=4,\qquad
d_6^{\mathrm{bin}}(e_3)=d_6^{\mathrm{bin}}(-e_3)=4.
\]
因此，window-\(6\) 的 bare fold potential 在长根层完全射影化，而在短根层仅于 \(\{\pm e_1\}\) 上出现唯一的定向破缺。
\end{corollary}
\begin{proof}
由定理 \ref{thm:conclusion-window6-coxeter-three-necklace-decomposition}，\(\mathfrak c^3\) 对应于根空间中的 \(-I\)，故每条 \(6\)-循环中相隔三步的两点恰为一对 antipodal roots。
再由定理 \ref{thm:conclusion-window6-coxeter-fold-block-normal-form}，
\[
D_{\ell}=\operatorname{diag}(4,2,3,4,2,3),
\]
前三项与后三项逐一相同，故长根上满足
\[
d_6^{\mathrm{bin}}(\alpha)=d_6^{\mathrm{bin}}(-\alpha).
\]
同样，
\[
D_{\mathrm{s}}=\operatorname{diag}(4,4,4,2,4,4)
\]
只在第一对 antipodal 位置上发生失配，而按照 \(\mathcal O_{\mathrm{s}}\) 的顺序，这一对恰对应
\[
\texttt{100000}\leftrightarrow e_1,\qquad
\texttt{000001}\leftrightarrow -e_1.
\]
于是短根层的射影对称只在 \(\{\pm e_1\}\) 上失败。证毕。
\end{proof}

\begin{theorem}[长根项链的 Fourier 选择律]\label{thm:conclusion-window6-long-root-fourier-selection-rule}
\leanverified{paper\_conclusion\_window6\_long\_root\_fourier\_selection\_rule}
在定理 \ref{thm:conclusion-window6-coxeter-three-necklace-decomposition} 的循环顺序下，长根块的 multiplicity profile 为
\[
p=(4,2,3,4,2,3).
\]
定义其 \(6\)-点离散 Fourier 系数
\[
\widehat p(k):=\sum_{j=0}^{5}p_j\,\zeta_6^{-kj},
\qquad
\zeta_6=e^{2\pi i/6}.
\]
则
\[
\widehat p(1)=\widehat p(3)=\widehat p(5)=0.
\]
等价地，长根块 \(D_{\ell}\) 的 orbitwise Fourier 展开仅支撑于偶 Coxeter 动量 \(k=0,2,4\)。

相反，短根块的 profile
\[
q=(4,4,4,2,4,4)
\]
满足
\[
\widehat q(k)=24\delta_{k0}-2(-1)^k,
\]
其中 \(\delta_{k0}\) 为 Kronecker 符号。故短根项链的六个 Fourier 模全部非零。
\end{theorem}
\begin{proof}
对长根 profile \(p\)，有
\[
p_{j+3}=p_j\qquad (0\le j\le 2).
\]
因此
\[
\widehat p(k)
=
\sum_{j=0}^{2}p_j\,\zeta_6^{-kj}\bigl(1+\zeta_6^{-3k}\bigr).
\]
若 \(k\) 为奇数，则 \(\zeta_6^{-3k}=-1\)，故
\[
\widehat p(k)=0
\qquad (k=1,3,5).
\]
这说明长根块只含偶 Coxeter 动量。

对短根 profile \(q\)，可写成
\[
q_j=4-2\delta_{j3}
\qquad (0\le j\le 5).
\]
于是
\[
\widehat q(k)
=
4\sum_{j=0}^{5}\zeta_6^{-kj}-2\zeta_6^{-3k}
=
24\delta_{k0}-2(-1)^k.
\]
当 \(k=0\) 时，\(\widehat q(0)=22\)；当 \(1\le k\le 5\) 时，
\[
\widehat q(k)=-2(-1)^k\neq 0.
\]
故短根项链的六个 Fourier 模无一消失。证毕。
\end{proof}

\begin{theorem}[双长根项链的精确内部 \texorpdfstring{$GL_2$}{GL2} 对称性]\label{thm:conclusion-window6-long-root-gl2-symmetry}
令
\[
A_{\ell}:=\CC^{\mathcal O_{+}}\oplus \CC^{\mathcal O_{-}}.
\]
则存在规范同构
\[
A_{\ell}\cong \CC^6\otimes \CC^2_{\mathrm{orb}}
\]
使得
\[
U_{\mathfrak c}|_{A_{\ell}}=P_6\otimes I_2,\qquad
D_6|_{A_{\ell}}=D_{\ell}\otimes I_2.
\]
因此
\[
I_6\otimes M_2(\CC)\subseteq
\left\{
T\in \operatorname{End}(A_{\ell}):
\ TU_{\mathfrak c}=U_{\mathfrak c}T,\ TD_6=D_6T
\right\}.
\]
换言之，长根层不仅带有 Coxeter 六循环结构，而且带有一个与 bare fold potential 完全相容的精确内部 \(GL_2\) 对称；在酉归一化下，其紧实部分即给出 \(U(2)\) 对称，并进而产生 \(SU(2)\) doublet 结构。
\end{theorem}
\begin{proof}
由定理 \ref{thm:conclusion-window6-coxeter-fold-block-normal-form}，两条长根项链上的 Coxeter 块完全相同，均为 \(P_6\)；其 multiplicity 块也完全相同，均为 \(D_{\ell}\)。令
\[
e_j^{(+)}\in \CC^{\mathcal O_{+}},\qquad
e_j^{(-)}\in \CC^{\mathcal O_{-}}
\qquad (0\le j\le 5)
\]
分别为按循环顺序排列的标准基，并取 \(\CC^2_{\mathrm{orb}}\) 的标准基 \(f_{+},f_{-}\)。映射
\[
e_j^{(+)}\longmapsto e_j\otimes f_{+},\qquad
e_j^{(-)}\longmapsto e_j\otimes f_{-}
\]
给出一个规范同构
\[
A_{\ell}\cong \CC^6\otimes \CC^2_{\mathrm{orb}}.
\]
在该同构下，
\[
U_{\mathfrak c}|_{A_{\ell}}=P_6\otimes I_2,\qquad
D_6|_{A_{\ell}}=D_{\ell}\otimes I_2.
\]
于是任意形如 \(I_6\otimes T\) 的算子都同时与这两个张量积算子对易，故
\[
I_6\otimes M_2(\CC)\subseteq
\left\{
T\in \operatorname{End}(A_{\ell}):
\ TU_{\mathfrak c}=U_{\mathfrak c}T,\ TD_6=D_6T
\right\}.
\]
证毕。
\end{proof}

\begin{corollary}[非交换 Hamiltonian 的偶简并]\label{cor:conclusion-window6-even-degeneracy-noncommutative-hamiltonian}
设 \(H\) 是由
\[
U_{\mathfrak c}|_{A_{\ell}},\qquad
U_{\mathfrak c}^{-1}|_{A_{\ell}},\qquad
D_6|_{A_{\ell}}
\]
经任意非交换多项式构造得到的算子，即
\[
H=P\bigl(U_{\mathfrak c}|_{A_{\ell}},U_{\mathfrak c}^{-1}|_{A_{\ell}},D_6|_{A_{\ell}}\bigr)
\]
其中 \(P\in \CC\langle X,Y,Z\rangle\)。则存在某个 \(6\times 6\) 矩阵 \(H_6\) 使
\[
H=H_6\otimes I_2.
\]
因此 \(H\) 的每个特征值都具有偶数代数重数。特别地，任何仅由 Coxeter 传输与 bare fold potential 生成的 long-sector 动力学，都不可能在内部打破 doublet 简并。
\end{corollary}
\begin{proof}
由定理 \ref{thm:conclusion-window6-long-root-gl2-symmetry}，
\[
U_{\mathfrak c}|_{A_{\ell}}=P_6\otimes I_2,\qquad
U_{\mathfrak c}^{-1}|_{A_{\ell}}=P_6^{-1}\otimes I_2,\qquad
D_6|_{A_{\ell}}=D_{\ell}\otimes I_2.
\]
对任意非交换多项式 \(P\)，第二因子始终保持为 \(I_2\)，故
\[
H=P(P_6,P_6^{-1},D_{\ell})\otimes I_2=:H_6\otimes I_2.
\]
于是
\[
\chi_H(\lambda)=\chi_{H_6}(\lambda)^2,
\]
故 \(H\) 的每个特征值都具有偶数代数重数。证毕。
\end{proof}

\begin{theorem}[长根层的 qutrit--sign--doublet 张量正规形]\label{thm:conclusion-window6-qutrit-sign-doublet-normal-form}
长根层 \(A_{\ell}\) 还可进一步规范分解为
\[
A_{\ell}\cong \CC^3\otimes \CC^2_{\mathrm{sgn}}\otimes \CC^2_{\mathrm{orb}},
\]
其中第一因子记录每条 long Coxeter orbit 内的 projective 位置，第二因子记录 antipodal sign，第三因子记录两条 long Coxeter orbit 的标签。相对于适当基底，有
\[
D_6|_{A_{\ell}}=\operatorname{diag}(4,2,3)\otimes I_2\otimes I_2,
\]
并且
\[
U_{\mathfrak c}|_{A_{\ell}}=Q_3\otimes S_2\otimes I_2,
\]
其中 \(Q_3\) 为标准 \(3\)-循环矩阵，\(S_2\) 为 \(2\times 2\) 交换矩阵。

因此，window-\(6\) 的长根层精确分裂为一个承载三值 fold 势 \((4,2,3)\) 的 projective qutrit、一个记录 \(c^3=-I\) 的 sign 因子，以及一个记录双长根项链内部交换对称的 orbit doublet。
\end{theorem}
\begin{proof}
先固定一条长根项链，并以 \(t\in \ZZ/6\ZZ\) 记录其循环位置。写
\[
a:=t\bmod 3\in \ZZ/3\ZZ,\qquad
b:=t\bmod 2\in \ZZ/2\ZZ.
\]
由于映射
\[
t\longmapsto (a,b)
\]
给出显式同构
\[
\ZZ/6\ZZ\cong \ZZ/3\ZZ\times \ZZ/2\ZZ,
\]
故有规范识别
\[
\CC^6\cong \CC^3\otimes \CC^2_{\mathrm{sgn}}.
\]
再附加两条长根项链的 orbit 标记，即得到
\[
A_{\ell}\cong \CC^3\otimes \CC^2_{\mathrm{sgn}}\otimes \CC^2_{\mathrm{orb}}.
\]

由定理 \ref{thm:conclusion-window6-coxeter-fold-block-normal-form}，
\[
D_{\ell}=(4,2,3,4,2,3),
\]
故它只依赖于 \(a=t\bmod 3\)，从而
\[
D_6|_{A_{\ell}}=\operatorname{diag}(4,2,3)\otimes I_2\otimes I_2.
\]
另一方面，Coxeter 生成元在位置标号上作用为 \(t\mapsto t+1\)。在 \((a,b)\)-坐标下，这正变为
\[
(a,b)\longmapsto (a+1,b+1),
\]
故其矩阵恰为
\[
Q_3\otimes S_2.
\]
再考虑 orbit 标签不变，即得
\[
U_{\mathfrak c}|_{A_{\ell}}=Q_3\otimes S_2\otimes I_2.
\]
证毕。
\end{proof}

\paragraph{统一读法}
上述链条表明，window-\(6\) 的 Lie--information 统一并不止于 \(\mathrm{Spin}(7)\) 与八点超立方体的层面。表 \ref{tab:fold6_b3c3_root_dictionary_b3} 上的 \(B_3\) Coxeter 元把 \(X_6^{\mathrm{cyc}}\) 精确组织为一条短根六项链与两条长根六项链；bare fold potential 在两条长根项链上给出同一半周期 profile
\[
(4,2,3,4,2,3),
\]
其 Fourier 支撑仅剩偶 Coxeter 动量，因而 \(12\) 维 long-sector 被迫分裂为
\[
\CC^3\otimes \CC^2_{\mathrm{sgn}}\otimes \CC^2_{\mathrm{orb}}
\]
的 qutrit--sign--doublet 正规形。由此，任何只由 Coxeter 传输与 bare fold potential 生成的 long-layer Hamiltonian 都必然保有 doublet 简并；而整个 window-\(6\) 多重度场对射影根线的唯一破缺，则被压缩到短根层的单一方向 \(\{\pm e_1\}\) 上。换言之，连续 Lie 型、离散 Coxeter 动力学与 fold 势读出在这里并非三条平行叙述，而是在同一 \(12\) 维长根层内闭合为一个可完全张量化的刚性系统。

\endinput
